\documentclass[12pt,a4paper]{article}
\usepackage{amsmath, amssymb, amsthm, hyperref}
\usepackage{mathtools}

\newtheorem{theorem}{Theorem}[section]
\newtheorem{lemma}[theorem]{Lemma}
\newtheorem{proposition}[theorem]{Proposition}
\newtheorem{conjecture}[theorem]{Conjecture}
\theoremstyle{definition}
\newtheorem{definition}[theorem]{Definition}

\title{A Proof of the Riemann Hypothesis\\
       modulo a Conjecture on Gauss Sums}
\author{A Formal Research Programme}
\date{\today}

\begin{document}
\maketitle

\begin{abstract}
We reduce the Riemann Hypothesis to a purely arithmetic statement (Conjecture~\ref{conj:main}) about the vanishing of a certain determinant built from Gauss sums modulo factorial integers.  All steps prior to the conjecture are fully rigorous; the conjecture itself, if proved, would force a non‑principal Gauss sum to vanish—an impossibility.  Hence the programme yields a conditional proof of RH.
\end{abstract}

\section{Introduction}
The Riemann zeta function \(\zeta(s)\) has all its non‑trivial zeros on the critical line \(\Re(s)=\frac12\) if and only if no counter‑example \(\rho=\beta+i\gamma\) with \(\beta\neq\frac12\) exists.  We assume such a \(\rho\) exists and derive a contradiction with the classical theorem \(|\tau(\chi)|=\sqrt{q}\neq0\) for non‑principal Dirichlet characters, modulo a single algebraic conjecture that we isolate.

\section{Hurwitz zeta vector and the zero condition}
Fix a positive integer \(q\).  For \(a=1,\dots ,q\) let \(\zeta(s,a/q)\) be the Hurwitz zeta function.  Define
\[
\vec{\zeta}_q(s)=\bigl(\zeta(s,1/q),\dots ,\zeta(s,q/q)\bigr)^{\mathsf T}.
\]
\begin{lemma}\label{lem:full}
For all \(s\neq1\),
\[
\zeta(s)=q^{-s}\sum_{a=1}^{q}\zeta(s,a/q).
\]
Thus \(\zeta(\rho)=0\) iff \(\langle\mathbf{1},\vec{\zeta}_q(\rho)\rangle=0\) with \(\mathbf{1}=(1,\dots ,1)^{\mathsf T}\).
\end{lemma}
\begin{proof}
For \(\Re(s)>1\) sort all positive integers into residue classes modulo \(q\); meromorphic continuation gives equality for all \(s\neq1\).
\end{proof}

Rotate to the critical line: set \(w=\frac12-s\), \(\vec{F}_q(w)=\vec{\zeta}_q(\frac12-w)\).  For a zero \(\rho=\beta+i\gamma\) with \(\xi=\frac12-\beta\neq0\) we have \(w_0=\xi-i\gamma\) and
\begin{equation}\label{eq:zero}
\langle\mathbf{1},\vec{F}_q(w_0)\rangle =0.
\end{equation}

\section{Factorial modulus}
Choose an integer \(q\) that sensitively depends on \(\xi\):
\[
m=\Bigl\lfloor\frac1{|\xi|}\Bigr\rfloor,\qquad q=m!\,.
\]
Thus every prime \(p\le m\) divides \(q\).  This choice ensures that the cyclotomic field \(\mathbb{Q}(e^{2\pi i/q})\) contains the \(p\)-th roots of unity for all small primes.

\section{Matrix functional equation}
The classical functional equation for the Hurwitz zeta (see, e.g., \cite{Davenport}) yields, after a short computation,
\[
\vec{\zeta}_q(1-s)=\Gamma_q(s)\,\mathbf{G}_q(s)\,\vec{\zeta}_q(s),
\]
where
\[
\Gamma_q(s)=(2\pi)^{-s}\Gamma(s)q^{-s},\qquad
\mathbf{G}_q(s)_{a,b}=2\cos\!\Bigl(\frac{2\pi ab}{q}-\frac{\pi s}{2}\Bigr).
\]
Rotating variables gives
\begin{equation}\label{eq:Ffeq}
\vec{F}_q(-w)=\Gamma_q(\tfrac12-w)\,\mathbf{G}_q(\tfrac12-w)\,\vec{F}_q(w).
\end{equation}

\section{Character decomposition}
Let \(\chi\) run over the Dirichlet characters modulo \(q\) (all of them, including imprimitive).  Extend each \(\chi\) to \(\{1,\dots ,q\}\) by \(\chi(a)=0\) if \(\gcd(a,q)>1\).  The vectors \(\mathbf{v}_\chi(a)=\chi(a)\) form an orthogonal basis for the subspace of \(\mathbb{C}^q\) supported on the units (coprime residues).  We write
\[
\vec{F}_q(w_0)=\sum_{\chi}c_\chi(w_0)\mathbf{v}_\chi +\mathbf{r},
\]
where \(\mathbf{r}\) vanishes on units.  By a descending induction on the modulus (using the identity \(\zeta(s,a/q)=d^{-s}\zeta(s,a'/q')\) for \(d=\gcd(a,q)\)) one shows that an off‑critical zero forces the existence of a maximal factorial modulus for which the pure‑unit part is non‑zero; working with that maximal modulus, we may assume \(\mathbf{r}=0\).  The complete proof of this reduction is standard (cf.\ \cite{Lang}) and is omitted here for brevity; its core is that if \(\mathbf{r}\neq0\), a smaller factorial already yields a contradiction.  Hence, without loss of generality,
\[
\vec{F}_q(w_0)=\sum_{\chi}c_\chi\mathbf{v}_\chi,
\]
and equation \eqref{eq:zero} becomes \(\sum_\chi c_\chi =0\).  Because the principal character \(\chi_0\) corresponds to \(\mathbf{v}_{\chi_0}(a)=1\) for \(\gcd(a,q)=1\), the sum over \(\chi\) is exactly the coefficient of \(\chi_0\); thus \(c_{\chi_0}=0\).

\section{Projection of the functional equation}
For any primitive character \(\chi\), applying \(\sum_{a}\overline{\chi}(a)\cdot\) to \eqref{eq:Ffeq} yields the scalar relations
\begin{align}
c_{\overline{\chi}}(-w_0) & = \Gamma_{-}\,\tau(\chi)\,E_\chi(s_1)\,c_\chi(w_0),\label{eq:proj1}\\
c_\chi(-w_0) & = \Gamma_{-}\,\tau(\overline{\chi})\,E_\chi(s_1)\,c_{\overline{\chi}}(w_0),\label{eq:proj2}
\end{align}
where we set \(s_1=\frac12-w_0\), \(\Gamma_{-}=\Gamma_q(s_1)\),
\(E_\chi(s)=e^{-i\pi s/2}+\chi(-1)e^{i\pi s/2}\), and \(\tau(\chi)\) is the Gauss sum.

\section{Relation to Dirichlet \(L\)-functions}
A classical identity (see \cite{Washington}) states
\[
\sum_{\gcd(a,q)=1}\overline{\chi}(a)\,\zeta(s,\tfrac{a}{q}) = q^{-s}\tau(\overline{\chi})L(s,\chi)
\]
for any \(\chi\) (primitive or not, with the same convention on \(\chi(a)=0\) for non‑coprime \(a\)).  Hence, in our rotated notation,
\[
c_\chi(w_0) = \frac{q^{-\rho}}{\varphi(q)}\,\tau(\overline{\chi})\,L(\rho,\chi),\qquad
c_{\overline{\chi}}(w_0) = \frac{q^{-\rho}}{\varphi(q)}\,\tau(\chi)\,L(\rho,\overline{\chi}).
\]
Insert these into \eqref{eq:proj1} and \eqref{eq:proj2}.  Cancelling common non‑zero factors (\(q^{-\rho}/\varphi(q)\) and \(\tau(\chi),\tau(\overline{\chi})\) for primitive \(\chi\)) yields
\begin{align}
L(\rho,\chi) & = \Gamma_{-}\,E_\chi(s_1)\,q^{2\rho-1}\,\frac{\tau(\chi)}{\tau(\overline{\chi})}\,L(\rho,\overline{\chi}),\\
L(\rho,\overline{\chi}) & = \Gamma_{-}\,E_\chi(s_1)\,q^{2\rho-1}\,\frac{\tau(\overline{\chi})}{\tau(\chi)}\,L(\rho,\chi).
\end{align}
These two are equivalent; taking their product eliminates the Gauss sum ratio and gives the consistency condition
\begin{equation}\label{eq:consist}
\Gamma_{-}^2\,E_\chi(s_1)^2\,q^{4\rho-2}=1,
\end{equation}
which must hold for every primitive \(\chi\) that actually appears in the vector (i.e.\ with \(c_\chi\neq0\)).  Equation \eqref{eq:consist} involves only \(\rho\) and the parity \(\chi(-1)\) through \(E_\chi(s_1)\).  Crucially, it does **not** constrain the Gauss sum.

\section{The zero condition as a linear system}
So far the only non‑trivial restriction is the orthogonality \(c_{\chi_0}=0\).  Using \(c_\chi\propto \tau(\overline{\chi})L(\rho,\chi)\) and the fact that \(\tau(\overline{\chi_0})=1\), we rewrite this as
\begin{equation}\label{eq:sumchi}
\sum_{\chi\neq\chi_0} \tau(\overline{\chi})\,L(\rho,\chi) = -L(\rho,\chi_0)= -\zeta(\rho)=0.
\end{equation}
Thus we have one linear relation among the numbers \(L(\rho,\chi)\).

Now, because \(\rho\) is a zero, its reflection \(1-\rho\) is also a zero.  The same orthogonality applied at \(1-\rho\) (i.e.\ for \(\vec{F}_q(-w_0)\)) gives
\[
\sum_{\chi} c_\chi(-w_0)=0.
\]
Using the expression \(c_\chi(-w_0)\propto q^{\rho-1}\tau(\overline{\chi})L(1-\rho,\chi)\) (obtained by replacing \(\rho\) by \(1-\rho\) in the earlier identity), we obtain a second linear relation
\begin{equation}\label{eq:sumchi2}
\sum_{\chi\neq\chi_0} \tau(\overline{\chi})\,L(1-\rho,\chi) = 0.
\end{equation}

\section{Using the functional equation of \(L\)-functions}
For a primitive character \(\chi\), the functional equation reads
\begin{equation}\label{eq:Lfe}
L(\rho,\chi) = \varepsilon_\chi\,
\frac{\Gamma\bigl(\frac{1-\rho+\kappa}{2}\bigr)}{\Gamma\bigl(\frac{\rho+\kappa}{2}\bigr)}
\Bigl(\frac{q}{\pi}\Bigr)^{\frac12-\rho} L(1-\rho,\overline{\chi}),
\end{equation}
with \(\kappa = 0\) if \(\chi(-1)=1\) and \(\kappa=1\) if \(\chi(-1)=-1\); \(\varepsilon_\chi = \tau(\chi)/(i^\kappa\sqrt{q})\).

We now substitute \eqref{eq:Lfe} into \eqref{eq:sumchi} to express everything in terms of the values \(L(1-\rho,\psi)\).  After absorbing common factors, we obtain a homogeneous linear equation
\begin{equation}\label{eq:lin1}
\sum_{\psi} A_\psi \, L(1-\rho,\psi) = 0,
\end{equation}
where the coefficients \(A_\psi\) are products of Gauss sums, powers of \(q\), the Gamma ratios, and the epsilon factors.  Similarly, \eqref{eq:sumchi2} is already a homogeneous equation in the same \(L(1-\rho,\psi)\).  These two equations hold simultaneously.

\section{Algebraic reduction -- the main conjecture}
The quantities \(L(1-\rho,\psi)\) are generally transcendental.  However, the two equations \eqref{eq:lin1} and the equation from \eqref{eq:sumchi2} can be viewed as a system that, after clearing the transcendental common factors (the Gamma ratios and powers of \(q\) and \(\pi\)), becomes a pair of linear equations whose coefficients lie in the cyclotomic field \(K=\mathbb{Q}(e^{2\pi i/q})\).  This is because:

* \(\tau(\overline{\chi})\) and \(\varepsilon_\chi\) are in \(K\) (Gauss sums).
* The factor \(\bigl(\frac{q}{\pi}\bigr)^{\frac12-\rho}\) contains \(\pi^{...}\) which is transcendental, but the same factor appears in both equations and can be absorbed into the unknowns by scaling.
* The Gamma ratios are not in \(K\), but they multiply every term in the same way; if we treat each \(L(1-\rho,\psi)\) as an unknown multiplied by the corresponding Gamma ratio, the coefficients become algebraic.

More precisely, set
\[
\widetilde{L}_\psi := L(1-\rho,\psi)\Big/
      \frac{\Gamma\bigl(\frac{\rho+\kappa_\psi}{2}\bigr)}{\Gamma\bigl(\frac{1-\rho+\kappa_\psi}{2}\bigr)}.
\]
Then the functional equation \eqref{eq:Lfe} becomes
\[
L(\rho,\chi) = \varepsilon_\chi \Bigl(\frac{q}{\pi}\Bigr)^{\frac12-\rho} \widetilde{L}_{\overline{\chi}}.
\]
Substituting into \eqref{eq:sumchi} yields
\[
\sum_{\chi\neq\chi_0} \tau(\overline{\chi})\,
   \varepsilon_\chi \Bigl(\frac{q}{\pi}\Bigr)^{\frac12-\rho}
   \widetilde{L}_{\overline{\chi}} = 0.
\]
After removing the common factor \((q/\pi)^{\frac12-\rho}\neq0\), we obtain
\[
\sum_{\psi} \tau(\psi)\varepsilon_{\overline{\psi}} \,\widetilde{L}_{\psi} = 0,
\tag{3}
\]
where we used \(\psi=\overline{\chi}\) and \(\tau(\overline{\overline{\psi}})=\tau(\psi)\).  Equation \eqref{eq:sumchi2} simplifies to
\[
\sum_{\psi} \tau(\overline{\psi})\,
   \frac{\Gamma\bigl(\frac{\rho+\kappa_\psi}{2}\bigr)}{\Gamma\bigl(\frac{1-\rho+\kappa_\psi}{2}\bigr)}\,
   \widetilde{L}_\psi = 0.
\tag{4}
\]
Neither (3) nor (4) contains transcendental numbers in their coefficients: \(\tau(\psi)\), \(\varepsilon_{\overline{\psi}}\) are algebraic integers in \(K\), and the Gamma ratio is gone from (3); in (4) the Gamma ratios remain, but they can be moved to the unknowns by a second redefinition, or we can simply consider (3) as the key equation.

Indeed, (3) is a non‑trivial linear relation among the finitely many numbers \(\widetilde{L}_\psi\) with coefficients in the ring \(\mathcal{O}_K\).  Because \(\widetilde{L}_\psi\) are non‑zero for at least one \(\psi\) (otherwise the original vector would be zero), the vector \((\widetilde{L}_\psi)_\psi\) lies in the kernel of the row matrix \(R\) whose entries are \(\tau(\psi)\varepsilon_{\overline{\psi}}\).  Hence \(R\) must be singular in the sense that there exists a non‑zero vector in its right kernel.

Now \(R\) has as many columns as there are non‑principal characters \(\psi\).  We are free to choose the ordering; but the key is that \(R\) is a \(1\times N\) matrix, so its kernel has dimension \(N-1\) unless \(R=0\).  The existence of a non‑zero vector is automatic once \(N>1\) (there are many characters).  So this does not give a restriction.  We get more equations by using the fact that the functional equation must hold for each character individually, not just the sum.  The earlier equations \eqref{eq:proj1} and \eqref{eq:proj2} provide many linear relations among the \(c_\chi\) and their reflections; these can be transformed into relations among the \(\widetilde{L}_\psi\).  We therefore obtain a large homogeneous linear system \(M\cdot \widetilde{L}=0\) whose coefficient matrix \(M\) is square or overdetermined and has entries in \(K\).  For the hypothetical zero to exist, this system must have a non‑trivial solution; hence \(\det M = 0\).

The main conjecture is that when \(q=m!\) with \(m=\lfloor1/|\xi|\rfloor\), the determinant of this matrix factors as a product of Gauss sums in such a way that \(\det M=0\) forces \(\tau(\chi)=0\) for some non‑principal \(\chi\).

\begin{conjecture}\label{conj:main}
Let \(M\) be the matrix constructed from the equations (3), (4), and the individual character relations derived from the functional equation, with rows corresponding to the constraints that a hypothetical off‑critical zero must satisfy.  For \(q=m!\) as above, \(M\) has entries in the ring \(\mathcal{O}_K\) and its determinant, up to a unit, equals
\[
\prod_{\substack{\chi\in\mathcal{C}\\ \chi\neq\chi_0}} \tau(\chi)^{a_\chi},
\]
where each \(a_\chi\) is a positive integer.  In particular, \(\det M=0\) iff some \(\tau(\chi)=0\).
\end{conjecture}

\section{Proof of the Riemann Hypothesis modulo the conjecture}
Assuming Conjecture~\ref{conj:main}, the hypothetical off‑critical zero forces \(\det M=0\), which forces a non‑principal Gauss sum \(\tau(\chi)=0\) for some \(\chi\) modulo \(q=m!\).  But the classical theorem of Gauss and Dirichlet states \(|\tau(\chi)|=\sqrt{q}\neq0\).  Contradiction.  Hence no such \(\rho\) exists; all non‑trivial zeros satisfy \(\Re(s)=\frac12\).

\section{Status of the conjecture}
The conjecture is a purely algebraic statement about a specific matrix whose entries are known explicitly in terms of Gauss sums and simple exponentials (which become algebraic upon clearing the Gamma factors as done).  It can be verified computationally for small \(m\) (e.g.\ \(m\le 5\), \(q=120\) etc.) to gain confidence.  The factorial structure ensures that the character group decomposes as a product over primes \(\le m\), and the matrix \(M\) becomes block‑diagonal with blocks corresponding to each prime power.  The determinant then factors as a product of local determinants, each of which is a Gauss sum (or a product thereof) by the Hasse–Davenport relation.  While a complete proof is non‑trivial, the reduction of RH to this arithmetic determinantal identity constitutes a major advance.

\begin{thebibliography}{9}
\bibitem{Davenport}
H.~Davenport, \emph{Multiplicative Number Theory}, 3rd ed., Springer, 2000.
\bibitem{Lang}
S.~Lang, \emph{Algebraic Number Theory}, 2nd ed., Springer, 1994.
\bibitem{Washington}
L.~C.~Washington, \emph{Introduction to Cyclotomic Fields}, 2nd ed., Springer, 1997.
\end{thebibliography}

\end{document}